% B02-S001 | APPM5720Notes.tex baris 2750-2769
% segment-id: d90.becker.98ed693.b02.seg0001
\chapter{Pemisahan Douglas--Rachford}
\label{becker:chapter:douglas-rachford}

Bab ini menyajikan metode Douglas--Rachford untuk meminimumkan jumlah dua
fungsi konveks melalui operator proksimal. Notasi subdiferensial dan bentuk
dual dinormalkan agar persamaan konsisten dengan parameter skala yang dipakai
dalam algoritma.

\noindent\textbf{Kredit catatan donor.}\par
Catatan donor menyatakan bahwa bagian ini mengikuti Bauschke dan Combettes
(B\&C), edisi kedua (2017), \S20.3, dan mengatribusi analisis
Douglas--Rachford kepada Lions dan Mercier (1979). Kredit turunan tersebut
dipertahankan; edisi ini tidak mengimpor isi dari bagian yang bersebelahan.

Ambil $f,g\in\Gamma_0(\R^n)$, yaitu fungsi konveks proper dan semikontinu
bawah. Seperti pada donor, andaikan terdapat titik pelana primal--dual dan
aturan jumlah subdiferensial
\begin{equation}
  \label{becker:eq:dr-sum-rule}
  \partial(f+g)=\partial f+\partial g
\end{equation}
berlaku. Masalah primal adalah
\begin{equation}
  \label{becker:eq:dr-primal}
  (P)\qquad p^*=\min_{x\in\R^n}\bigl\{f(x)+g(x)\bigr\}.
\end{equation}
Dengan konjugat Fenchel
$h^*(u)=\sup_x\{\langle u,x\rangle-h(x)\}$, bentuk maksimisasi dual yang
konsisten adalah
\begin{equation}
  \label{becker:eq:dr-dual}
  (D)\qquad
  d^*=\sup_{u\in\R^n}\bigl\{-f^*(-u)-g^*(u)\bigr\}.
\end{equation}
Secara ekuivalen, titik dual yang sama meminimumkan
$f^*(-u)+g^*(u)$. Tanda pada bentuk dual donor diperbaiki secara terbuka
dalam edisi ini.

Untuk $h\in\Gamma_0(\R^n)$ dan $\rho>0$, tuliskan
\begin{equation}
  \label{becker:eq:dr-prox-definition}
  \prox_{\rho h}(v)
  =\operatorname*{arg\,min}_{w\in\R^n}
    \left\{h(w)+\frac{1}{2\rho}\norm{w-v}_2^2\right\}
  =(I+\rho\partial h)^{-1}(v).
\end{equation}
Pilih $y_0\in\R^n$, parameter skala $\rho>0$, dan parameter relaksasi tetap
$\lambda\in(0,2)$. Iterasi Douglas--Rachford adalah
\begin{equation}
  \label{becker:eq:dr-iteration}
  \begin{aligned}
    x_k&=\prox_{\rho g}(y_k),\\
    z_k&=\prox_{\rho f}(2x_k-y_k),\\
    y_{k+1}&=y_k+\lambda(z_k-x_k).
  \end{aligned}
\end{equation}
Catatan donor memperingatkan bahwa konvensi untuk $\rho$ pada bagian lain
mungkin memakai kebalikannya. Di sini $\rho$ selalu memiliki arti yang
ditetapkan oleh~\eqref{becker:eq:dr-prox-definition}.

% B02-S002 | APPM5720Notes.tex baris 2771-2790
% segment-id: d90.becker.98ed693.b02.seg0002
\section{Limit bayangan dan persamaan titik tetap}
\label{becker:section:dr-fixed-point}

\begin{theorem}[Limit bayangan merupakan solusi primal]
  \label{becker:theorem:dr-shadow-limit}
  Jika barisan pada~\eqref{becker:eq:dr-iteration} memenuhi
  $y_k\to\bar y$ dan
  \begin{equation}
    \bar x=\prox_{\rho g}(\bar y),
  \end{equation}
  maka $\bar x$ adalah solusi optimal
  masalah~\eqref{becker:eq:dr-primal}.
\end{theorem}
\begin{proof}
  Karena $y_k$ konvergen,
  $y_{k+1}-y_k=\lambda(z_k-x_k)\to0$, sehingga $z_k-x_k\to0$.
  Operator proksimal kontinu, maka
  \begin{equation}
    \bar x=\prox_{\rho g}(\bar y),
    \qquad
    \bar x=\prox_{\rho f}(2\bar x-\bar y).
  \end{equation}
  Syarat optimalitas kedua operator proksimal memberi
  \begin{equation}
    \frac{\bar y-\bar x}{\rho}\in\partial g(\bar x),
    \qquad
    \frac{\bar x-\bar y}{\rho}\in\partial f(\bar x).
  \end{equation}
  Menjumlahkan kedua inklusi dan memakai
  aturan~\eqref{becker:eq:dr-sum-rule} menghasilkan
  $0\in\partial(f+g)(\bar x)$, yang setara dengan optimalitas primal.
\end{proof}

Persamaan titik tetapnya dapat dilihat langsung dari syarat optimalitas.
Untuk solusi primal $x$, pilih $u\in\partial g(x)$ dengan
$-u\in\partial f(x)$, lalu definisikan $y=x+\rho u$. Dengan demikian
\begin{equation}
  \label{becker:eq:dr-resolvent-motivation}
  \begin{aligned}
    y&\in x+\rho\partial g(x),
      &x&=(I+\rho\partial g)^{-1}(y)=\prox_{\rho g}(y),\\
    2x-y=x-\rho u&\in x+\rho\partial f(x),
      &x&=(I+\rho\partial f)^{-1}(2x-y)=\prox_{\rho f}(2x-y).
  \end{aligned}
\end{equation}
Jadi, dengan
\begin{equation}
  x=\prox_{\rho g}(y),
  \qquad z=\prox_{\rho f}(2x-y),
\end{equation}
syarat solusi menjadi $z-x=0$. Pembaruan
$y\leftarrow y+\lambda(z-x)$ pada~\eqref{becker:eq:dr-iteration}
merupakan iterasi relaksasi menuju titik tetap itu. Faktor $\rho$ dan pilihan
subgradien $u$ ditulis eksplisit karena subdiferensial pada umumnya bernilai
himpunan dan tidak dapat diperlakukan sebagai satu vektor $d g(x)$.

% B02-S003 | APPM5720Notes.tex baris 2792-2797
% segment-id: d90.becker.98ed693.b02.seg0003
\section{Hubungan dengan ADMM}
\label{becker:section:dr-admm}

Pernyataan singkat donor bahwa ADMM merupakan kasus khusus
Douglas--Rachford dipahami dalam bentuk yang lazim dan lebih presisi: ADMM
dapat diperoleh sebagai pemisahan Douglas--Rachford yang diterapkan pada
masalah dual yang sesuai. Bagian ADMM yang mendahului rentang beku tidak
diimpor ke modul ini.
